\section{Computational expenditure as part of the design}

\subsection{Specify the workload before comparing methods}

A time comparison is meaningful only after defining the computational unit. For a collection of frozen policies under one financial model, approximate-value construction and the common martingale can be shared. Charging this construction once per policy overstates its cost; omitting it understates the cost. A single-policy workload can have a different ranking from a repeated multi-policy workload.

The timed unit should state whether it includes simulation, deterministic approximation, conditional integration, martingale evaluation, policy realisation, estimator construction and numerical checks. Previous training, imports, data loading and file output should be reported separately if excluded. Warm-up and allocation pilots are additional expenses, not free inputs to an end-to-end valuation. Record elapsed and CPU time and the relevant hardware and thread configuration.

\subsection{Allocate paths before inspecting evaluation outcomes}

One possible allocation design uses median pilot times $t_1,t_2$ at two path counts $n_1<n_2$:
\begin{equation}
 b_T=\frac{t_2-t_1}{n_2-n_1},\qquad
 a_T=\max(0,t_1-b_Tn_1),\qquad
 \widehat T(n)=a_T+b_Tn.
\label{eq:timing}
\end{equation}
For target $B$ and path increment $q$, choose
\begin{equation}
 n(B)=q\left\lfloor\frac{B-a_T}{b_Tq}\right\rfloor,
\label{eq:allocation}
\end{equation}
subject to prescribed limits. A non-positive fitted slope or infeasible minimum is a design failure. The unclipped intercept, timing dispersion and any cap must be retained. Fix all allocations before drawing the independent evaluation streams, and do not stop a run because a confidence bound becomes favourable.

The affine model predicts time; it does not enforce it. Extrapolation beyond the pilot range, memory effects and host variability can create substantial discrepancies. Compare actual time with target time and retain overruns. A strict equal-cost ranking requires adequate timing calibration or a separately justified comparison on an observed cost frontier. Nominally equal targets alone do not establish it.

\subsection{Approximation and sampling serve different purposes}

A more detailed approximate value can tighten the dual expectation while making each outer path more expensive. Under a time constraint this can reduce the available sample size and increase the sampling radius in \eqref{eq:pointwise}. Report the mean gap, its standard error, the upper limit and actual runtime separately. A finer deterministic grid is not automatically the best use of a limited evaluation budget.

For independent-method evaluations, use \eqref{eq:welch}; for common-path evaluations, retain \eqref{eq:paired}. Report all retained policies and repetitions. Descriptive averages or medians across related cases are not independent observations of a population effect. Neither a tighter dual construction nor a larger outer sample changes a frozen policy's expected reward. Comparing these expenditures with extra policy training is a distinct resource-allocation problem.
